% Poincaré duality: cap with [M] gives k <-> n-k.
\begin{tikzpicture}[
  every node/.style={font=\small},
  arrow/.style={-{Stealth[length=2.5mm]}, thick},
]
  % Closed manifold M (drawn as a torus to remind dimension n)
  \draw[blue!60!black, thick, fill=blue!8] (0, 0) ellipse [x radius=2.0cm, y radius=1.0cm];
  \draw[blue!60!black, thick] (-0.6, 0.05) .. controls (-0.4, -0.25) and (0.4, -0.25) .. (0.6, 0.05);
  \draw[blue!60!black, thick] (-0.5, -0.05) .. controls (-0.3, 0.1) and (0.3, 0.1) .. (0.5, -0.05);

  \node[blue!60!black] at (-2.4, 0.7) {$M^n$};
  \node[font=\footnotesize, gray!50!black] at (0, -1.5) {closed oriented};

  % Cap-with-[M] arrow
  \draw[arrow] (2.5, 0) -- node[above] {$- \frown [M]$} node[below, font=\footnotesize] {Poincaré duality} (5.5, 0);

  % H^k -> H_{n-k}
  \node (Hk)  at (3.0, 1.6) {$H^k(M; \mathbb{Z})$};
  \node (Hnk) at (5.7, 1.6) {$H_{n-k}(M; \mathbb{Z})$};
  \draw[arrow] (Hk) -- node[above] {$\sim$} (Hnk);

  % Side: cap product table
  \node[anchor=west] at (6.5, -0.3) {Effects:};
  \node[anchor=west, font=\footnotesize] at (6.5, -0.9) {$H^0 \cong H_n$ (orientation class)};
  \node[anchor=west, font=\footnotesize] at (6.5, -1.4) {$H^k \cong H_{n-k}$};
  \node[anchor=west, font=\footnotesize] at (6.5, -1.9) {$H^n \cong H_0 = \mathbb{Z}$};

  \begin{scope}[on background layer]
    \fill[purple!4, rounded corners=12pt] (-3.2, -2.5) rectangle (10.5, 2.6);
  \end{scope}

  \node[align=center] at (3.0, -3.0)
    {\itshape For closed oriented $n$-manifolds, $\mathbb{Z}/2$-coefficients give the same iso without orientation.};
\end{tikzpicture}
